% GENERATED FILE. Edit canonical lessons, then run npm run generate:books.
% canonical-source-hash: fnv1a64:a54e5ee4544f0757
% canonical-lessons: LA-C308-iudex, LA-C308-iudicium, LA-C308-iudico, LA-C308-lex, LA-C308-ius

\chapter{Iūdex et lēx}
\label{ch:la-iudex-et-lex}
\hlchaptermodality{308}

\begin{chapteropening}
\textbf{By the end of this chapter:} \emph{I can talk about a Roman court: the judge, the trial, to judge, a law and rights.}

Complete the last lesson of chapter 308: I can talk about a Roman court: the judge, the trial, to judge, a law and rights.
\end{chapteropening}

\section[iūdex, iūdicis]{iūdex, iūdicis — a judge; a juror}

\label{lesson:LA-C308-iudex}



\begin{quote}
Before the new one: say the Latin for a conspiracy, then the Latin for a traitor.
\end{quote}

\subsection*{You'll want to know: iūdex, iūdicis}
\textbf{iūdex, iūdicis} — \textquotedblleft{}a judge; a juror\textquotedblright{}.

\begin{quote}
Iūdex testēs diū audīvit.
\end{quote}

\begin{quote}
\emph{The judge listened to the witnesses for a long time.}
\end{quote}

Masculine (also feminine), 3rd declension. In big trials many \emph{iūdicēs} sat together as a jury.

\begin{cousinweb}
\textbf{iūdex} joins \emph{iūs}, \textquotedblleft{}law, right\textquotedblright{}, and the root of \emph{dīcō}, \textquotedblleft{}to say\textquotedblright{} (EtymDB): the one who speaks the law. English \emph{judge} comes from it through French.
\end{cousinweb}

\begin{grammarlens}[title={What you've built}]
One more word for telling, explaining and describing at B1.
\end{grammarlens}

\subsection*{Your turn}
\begin{itemize}
\raggedright
  \item \emph{Say it:} \emph{iūdex, iūdicis}
  \item \emph{Say it:} \emph{iūdex, iūdicis}, once more
  \item \emph{Say it:} a sentence of your own with \emph{iūdex, iūdicis}
  \item \emph{Recall:} say the Latin for a conspiracy, then the Latin for a traitor, then say \emph{iūdex, iūdicis} again
\end{itemize}

\begin{tcolorbox}[breakable,colback=teal!4,colframe=teal!35!black,title={Before you move on}]
What does iūdex, iūdicis mean? (\textquotedblleft{}A judge; a juror\textquotedblright{}.) Say it once more.
\end{tcolorbox}
\section[iūdicium, iūdiciī]{iūdicium, iūdiciī — a trial; a judgement}

\label{lesson:LA-C308-iudicium}



\begin{quote}
Before the new one: say the Latin for a traitor, then the Latin for a judge.
\end{quote}

\subsection*{You'll want to know: iūdicium, iūdiciī}
\textbf{iūdicium, iūdiciī} — \textquotedblleft{}a trial; a judgement\textquotedblright{}.

\begin{quote}
Iūdicium in basilicā habitum est.
\end{quote}

\begin{quote}
\emph{The trial was held in the basilica.}
\end{quote}

Neuter, 2nd declension. Both the trial and the verdict at its end; also 'opinion': \emph{meō iūdiciō}, in my judgement.

\begin{cousinweb}
\textbf{iūdicium} is built on \emph{iūdex}, \textquotedblleft{}judge\textquotedblright{} (EtymDB): what a judge does, the trial and its verdict. English \emph{judicial} belongs to the same family.
\end{cousinweb}

\begin{grammarlens}[title={What you've built}]
One more word for telling, explaining and describing at B1.
\end{grammarlens}

\subsection*{Your turn}
\begin{itemize}
\raggedright
  \item \emph{Say it:} \emph{iūdicium, iūdiciī}
  \item \emph{Say it:} \emph{iūdicium, iūdiciī}, once more
  \item \emph{Say it:} a sentence of your own with \emph{iūdicium, iūdiciī}
  \item \emph{Recall:} say the Latin for a traitor, then the Latin for a judge, then say \emph{iūdicium, iūdiciī} again
\end{itemize}

\begin{tcolorbox}[breakable,colback=teal!4,colframe=teal!35!black,title={Before you move on}]
What does iūdicium, iūdiciī mean? (\textquotedblleft{}A trial; a judgement\textquotedblright{}.) Say it once more.
\end{tcolorbox}
\section[iūdicō, iūdicāre]{iūdicō, iūdicāre — to judge; to declare}

\label{lesson:LA-C308-iudico}



\begin{quote}
Before the new one: say the Latin for a judge, then the Latin for a trial.
\end{quote}

\subsection*{You'll want to know: iūdicō, iūdicāre}
\textbf{iūdicō, iūdicāre} — \textquotedblleft{}to judge; to declare\textquotedblright{}.

\begin{quote}
Senātus eum hostem iūdicāvit.
\end{quote}

\begin{quote}
\emph{The senate declared him a public enemy.}
\end{quote}

\emph{iūdicō, iūdicāre, iūdicāvī, iūdicātum}, often with two accusatives: someone + what he is judged to be.

\begin{cousinweb}
\textbf{iūdicō} is built on \emph{iūdex}, \textquotedblleft{}judge\textquotedblright{} (EtymDB): to act as judge, to decide. English \emph{judge} the verb comes from it through French.
\end{cousinweb}

\begin{grammarlens}[title={What you've built}]
One more word for telling, explaining and describing at B1.
\end{grammarlens}

\subsection*{Your turn}
\begin{itemize}
\raggedright
  \item \emph{Say it:} \emph{iūdicō, iūdicāre}
  \item \emph{Say it:} \emph{iūdicō, iūdicāre}, once more
  \item \emph{Say it:} a sentence of your own with \emph{iūdicō, iūdicāre}
  \item \emph{Recall:} say the Latin for a judge, then the Latin for a trial, then say \emph{iūdicō, iūdicāre} again
\end{itemize}

\begin{tcolorbox}[breakable,colback=teal!4,colframe=teal!35!black,title={Before you move on}]
What does iūdicō, iūdicāre mean? (\textquotedblleft{}To judge; to declare\textquotedblright{}.) Say it once more.
\end{tcolorbox}
\section[lēx, lēgis]{lēx, lēgis — a law}

\label{lesson:LA-C308-lex}



\begin{quote}
Before the new one: say the Latin for a trial, then the Latin for to judge.
\end{quote}

\subsection*{You'll want to know: lēx, lēgis}
\textbf{lēx, lēgis} — \textquotedblleft{}a law\textquotedblright{}.

\begin{quote}
Populus novam lēgem probāvit.
\end{quote}

\begin{quote}
\emph{The people approved the new law.}
\end{quote}

Feminine, 3rd declension. A \emph{lēx} was a bill passed by the people's vote. English \emph{legal} comes from it.

\begin{cousinweb}
\textbf{lēx} is inherited from Proto-Italic \emph{*lēg-} (EtymDB); Lewis \& Short connect it, with doubt, with a root meaning \textquotedblleft{}to bind\textquotedblright{}. English \emph{legal} and \emph{legislate} come from it.
\end{cousinweb}

\begin{grammarlens}[title={What you've built}]
One more word for telling, explaining and describing at B1.
\end{grammarlens}

\subsection*{Your turn}
\begin{itemize}
\raggedright
  \item \emph{Say it:} \emph{lēx, lēgis}
  \item \emph{Say it:} \emph{lēx, lēgis}, once more
  \item \emph{Say it:} a sentence of your own with \emph{lēx, lēgis}
  \item \emph{Recall:} say the Latin for a trial, then the Latin for to judge, then say \emph{lēx, lēgis} again
\end{itemize}

\begin{tcolorbox}[breakable,colback=teal!4,colframe=teal!35!black,title={Before you move on}]
What does lēx, lēgis mean? (\textquotedblleft{}A law\textquotedblright{}.) Say it once more.
\end{tcolorbox}
\section[iūs, iūris]{iūs, iūris — law, right; justice}

\label{lesson:LA-C308-ius}



\begin{quote}
Before the new one: say the Latin for to judge, then the Latin for a law.
\end{quote}

\subsection*{You'll want to know: iūs, iūris}
\textbf{iūs, iūris} — \textquotedblleft{}law, right; justice\textquotedblright{}.

\begin{quote}
Omnēs cīvēs eadem iūra habēbant.
\end{quote}

\begin{quote}
\emph{All the citizens had the same rights.}
\end{quote}

Neuter, 3rd declension. \emph{Iūs} is law and right as a whole; a \emph{lēx} is one law. A look-alike \emph{iūs} means 'broth'.

\begin{cousinweb}
\textbf{iūs} is inherited from Proto-Indo-European \emph{*h\textsubscript{2}yéwos}, \textquotedblleft{}law, right\textquotedblright{} (EtymDB). English \emph{jury}, \emph{justice} and \emph{jurist} come from its family.
\end{cousinweb}

\begin{grammarlens}[title={What you've built}]
One more word for telling, explaining and describing at B1.
\end{grammarlens}

\subsection*{Your turn}
\begin{itemize}
\raggedright
  \item \emph{Say it:} \emph{iūs, iūris}
  \item \emph{Say it:} \emph{iūs, iūris}, once more
  \item \emph{Say it:} a sentence of your own with \emph{iūs, iūris}
  \item \emph{Recall:} say the Latin for to judge, then the Latin for a law, then say \emph{iūs, iūris} again
\end{itemize}

\begin{tcolorbox}[breakable,colback=teal!4,colframe=teal!35!black,title={Before you move on}]
What does iūs, iūris mean? (\textquotedblleft{}Law, right; justice\textquotedblright{}.) Say it once more.
\end{tcolorbox}
